\section{Routley--Meyer semantics for relevance}
\label{sec:relevance-linear-02}

The formula language used in this section is exactly
\[
  A::=p\mid\neg A\mid(A\land A)\mid(A\lor A)\mid(A\to A),
  \qquad p\in\Prop.
\]
No consequence relation or proof calculus is determined by this formation
grammar alone.

\begin{definition}[Routley--Meyer frame and structure]
A Routley--Meyer frame is a tuple
\[
  \mathcal F=(W,N,R,{*}),
\]
where \(W\ne\varnothing\), \(N\subseteq W\) is a nonempty set of normal worlds,
\(R\subseteq W^3\), and \({*}:W\to W\) is an involution.  The base normality
condition used here is
\[
  n\in N
  \quad\Longrightarrow\quad
  (R(n,u,v)\Longleftrightarrow u=v).
\]
A Routley--Meyer structure is
\(\mathcal M=(W,N,R,{*},V)\), where the frame components satisfy the displayed
conditions and \(V:\Prop\to\Pow(W)\) is a valuation.
\end{definition}

Atomic satisfaction is
\[
  \mathcal M,w\entails p
  \quad\Longleftrightarrow\quad
  w\in V(p).
\]
Satisfaction for conjunction and disjunction is local and classical.  The
relevance-sensitive clauses are
\[
  \mathcal M,w\entails\neg A
  \quad\Longleftrightarrow\quad
  \mathcal M,w^*\notentails A,
\]
and
\[
  \mathcal M,w\entails A\to B
  \quad\Longleftrightarrow\quad
  \forall u,v\,
  \bigl((R(w,u,v)\land\mathcal M,u\entails A)
        \to\mathcal M,v\entails B\bigr).
\]
Validity requires satisfaction at every normal world.  Different relevance
logics impose additional first-order conditions on \(R\), \({*}\), and \(N\),
and use matching Hilbert, display, or labelled calculi.  No such calculus is
selected by the semantic clauses alone.

